\section*{Hoofdstuk \ref{ch:b2:elements-in-the-earth} --- De elementen in de aarde}

\begin{solution}{exo:b2:elements-in-the-earth:1}
Kalium is $2\times39.10/(2\times39.10+16.00) = 78.20/94.20 = 0.8301$ van de massa van
\ce{K2O}, dus $\num{5130}\times0.8301 = \qty{4259}{ppm}$ K. Titaan is
$47.87/79.87 = 0.5993$ van \ce{TiO2}: $\num{27310}\times0.5993 = \qty{16370}{ppm}$ Ti.
\end{solution}

\begin{solution}{exo:b2:elements-in-the-earth:2}
Aluminium is lithofiel en refractair: de silicaataarde heeft het helemaal behouden, in bijna
chondritische verhoudingen (1.16). Natrium is lithofiel maar vluchtig (\qty{958}{K}): het ging
vóór of tijdens de accretie verloren, en er bleef 0.22 over. Nikkel condenseert met magnesium
(\qty{1353}{K}), dus vluchtigheid kan 0.077 niet verklaren: het is siderofiel, en het grootste deel
ervan ging naar de kern. Iridium is refractair (\qty{1603}{K}) en sterk siderofiel:
0.0029, bijna alles ervan zit in de kern.
\end{solution}

\begin{solution}{exo:b2:elements-in-the-earth:3}
$\mathrm{La}_N = 31/0.232 = 133.6$, $\mathrm{Yb}_N = 2.0/0.163 = 12.27$, en
$(\mathrm{La}/\mathrm{Yb})_N = 10.9$. De korst is verrijkt in alle \omterm{def:b2:elements-in-the-earth:ree}{zeldzame aarden}, en
elf keer meer in de lichte dan in de zware: ze bestaat uit smelten, die
de incompatibele lichte \omterm{def:b2:elements-in-the-earth:ree}{zeldzame aarden} het gemakkelijkst opnemen.
\end{solution}

\begin{solution}{exo:b2:elements-in-the-earth:4}
$\mathrm{Sm}_N = 4.7/0.145 = 32.4$, $\mathrm{Eu}_N = 1.0/0.0546 = 18.3$,
$\mathrm{Gd}_N = 4.0/0.198 = 20.2$; $\mathrm{Eu}^* = \sqrt{32.4\times20.2} = 25.6$ en
$\mathrm{Eu}/\mathrm{Eu}^* = 18.3/25.6 = 0.72$, een negatieve anomalie. Plagioklaas, dat
\ce{Eu^2+} opneemt, kristalliseerde uit de smelten die de bovenste korst opbouwden en bleef
dieper achter.
\end{solution}

\begin{solution}{exo:b2:elements-in-the-earth:5}
$\bar D_{\mathrm{La}} = 0.65\times0.058+0.18\times0.001+0.17\times0.001 = 0.0377 +
0.0002 + 0.0002 = 0.038$; $\bar D_{\mathrm{Lu}} = 0.3445+1.044+0.0306 = 1.42$. Lutetium is
compatibel ($\bar D > 1$), lanthaan sterk incompatibel. Granaat draagt
$1.044/1.42 = 74$ procent van het lutetium van het gesteente, hoewel het maar 18 procent van het
gesteente uitmaakt.
\end{solution}

\begin{solution}{exo:b2:elements-in-the-earth:6}
Voor natrium is Na/Mg van de silicaataarde $2670/\num{228000} = 0.01171$ en de CI-waarde
$5010/\num{95870} = 0.05226$: de verhouding is 0.224. Voor nikkel $1960/\num{228000}
= 0.008596$ tegen $\num{10640}/\num{95870} = 0.1110$: 0.0775. Natrium is
verarmd door zijn vluchtigheid (\qty{958}{K}); nikkel condenseert boven \qty{1300}{K} en is
drie keer sterker verarmd dan natrium omdat het siderofiel is: de kern nam het grootste deel
ervan op.
\end{solution}

\begin{solution}{exo:b2:elements-in-the-earth:7}
La: $5.21\times0.058 = \qty{0.302}{ppm}$; Lu: $0.53\times0.53 = \qty{0.281}{ppm}$. De
normalisatiewaarden vallen weg in de verhouding van de verhoudingen:
$(\mathrm{La}/\mathrm{Lu})_{N,\mathrm{cpx}}/(\mathrm{La}/\mathrm{Lu})_{N,\mathrm{melt}} =
D_{\mathrm{La}}/D_{\mathrm{Lu}} = 0.058/0.53 = 0.11$: het patroon van het kristal is negen
keer steiler naar de zware \omterm{def:b2:elements-in-the-earth:ree}{zeldzame aarden} toe dan dat van de smelt.
\end{solution}

\begin{solution}{exo:b2:elements-in-the-earth:8}
Korst: $(12/0.658)/(31/0.648) = 18.24/47.84 = 0.38$. Hawaï:
$(18.10/0.658)/(15.20/0.648) = 27.51/23.46 = 1.17$. MORB: $7.96/8.04 = 0.99$. Alleen de
korst, met een sterk tekort aan niobium, draagt het kenmerk van subductiezones.
\end{solution}

\begin{solution}{exo:b2:elements-in-the-earth:9}
\omterm{def:b2:elements-in-the-earth:bse}{Primitieve mantel}: $240/0.0203 = \num{11800}$. MORB: K $= 0.160\times0.830\times10^4 =
\qty{1328}{ppm}$, K/U $= \num{11200}$. Korst: K $= 2.80\times0.830\times10^4 =
\qty{23240}{ppm}$, K/U $= \num{8600}$. Kalium stijgt honderdvoudig van mantel tot korst,
maar K/U verandert met minder dan een derde: de twee elementen zijn ongeveer even incompatibel,
en smelten scheidt ze nauwelijks. Uranium is refractair, dus zijn abundantie volgt uit
het chondritische model; K/U gemeten aan lava's geeft dan het kalium van de planeet, dat,
omdat het vluchtig is, niet uit \omterm{def:b2:planet-formation-meteorites:chondrule}{chondrieten} kan worden afgeleid.
\end{solution}

\begin{solution}{exo:b2:elements-in-the-earth:10}
In evenwicht geldt $C_{\mathrm{rock}} = \bar D\,C_{\mathrm{melt}}$ voor elk element (en een
kleine smeltfractie laat de samenstelling van het gesteente bijna ongewijzigd), dus
$(C_1/C_2)_{\mathrm{melt}} = (C_{1}/\bar D_1)/(C_{2}/\bar D_2)$ voor de $C_1,
C_2$ van het gesteente, en de verhouding is $\bar D_2/\bar D_1$. Voor La en Lu: $1.419/0.03805 = 37$. De smelt
van een granaathoudende mantel is 37 keer rijker aan La ten opzichte van Lu dan zijn bron: granaat
hield de zware \omterm{def:b2:elements-in-the-earth:ree}{zeldzame aarden} vast, en het patroon van de smelt is steil.
\end{solution}

\begin{solution}{exo:b2:elements-in-the-earth:11}
$C_{\mathrm{mix}} = fC_U + (1-f)C_M$ voor elk element: lineair in $f$. De verhouding
$\bigl(f\,\mathrm{Th}_U + (1-f)\mathrm{Th}_M\bigr)/\bigl(f\,\mathrm{Nb}_U +
(1-f)\mathrm{Nb}_M\bigr)$ is een quotiënt van twee lineaire functies, een hyperbool in $f$.
Th/Nb $= 0.2$ als $f(10.5-2.4) + (1-f)(0.404-1.048) = 0$, dat wil zeggen $8.1f = 0.644(1-f)$,
$f = 0.644/8.744 = 0.074$: ongeveer 7 procent korst verhoogt Th/Nb van 0.077 tot 0.2. Voor
$f = 0.1$: Th $= 1.05 + 0.364 = 1.414$, Nb $= 1.2 + 4.716 = 5.916$, Th/Nb $= 0.24$.
\end{solution}

\begin{solution}{exo:b2:elements-in-the-earth:12}
Het verwijderen van een component die geen van beide elementen bevat, laat beide massa's, en dus hun verhouding,
ongewijzigd. Als al het aluminium is behouden, betekent $(\mathrm{Al}/\mathrm{Mg}) = 1.16$ keer CI
dat de silicaataarde $1/1.16 = 0.86$ van haar chondritische magnesium bevat: 14
procent ontbreekt. Silicium ten opzichte van aluminium is $0.83/1.16 = 0.72$ keer de CI-waarde:
28 procent van het chondritische silicium zou ontbreken.
\end{solution}

\begin{solution}{pb:b2:elements-in-the-earth:1}
\textbf{1.} Fe/Mg $= \num{182800}/\num{95870} = 1.907$; Ni/Mg $= 0.1110$; Si/Mg $=
1.111$; Al/Mg $= 0.0887$.

\textbf{2.} $(\num{182800}/55.85)/(\num{106500}/28.09) = 3273/3791 = 0.863$ ijzeratoom per
siliciumatoom, binnen 7 procent van de 0.81 van de fotosfeer: zinvol, want \omterm{def:b2:elements-in-the-earth:ci}{CI-chondrieten}
komen voor de gesteentevormende elementen overeen met de Zon.

\textbf{3.} Het verlies van een component die geen van beide elementen van een verhouding bevat, laat de
verhouding ongewijzigd (\cref{exo:b2:elements-in-the-earth:12}): alleen de absolute concentraties veranderen.

\textbf{4.} Magnesium is lithofiel en gaat niet naar de kern; het is een hoofdelement,
goed bepaald in zowel de \omterm{def:b2:planet-formation-meteorites:chondrule}{chondrieten} als de mantel. IJzer is juist het element dat verdeeld wordt
tussen kern en mantel: het kan niet zijn eigen referentie zijn.

\textbf{5.} Elementen die samen condenseren, werden in de nevel niet door vluchtigheid
gescheiden, dus de aarde als geheel zou hun verhoudingen moeten behouden. Al/Mg wijkt af: aluminium is
refractair (\qty{1653}{K}), en Al/Mg van de silicaataarde is 1.16 keer de CI-waarde.

\textbf{6.} Fe $= 8.05\times55.85/71.85 = 6.26$ procent; Ni $= 0.25\times58.69/74.69 =
0.196$ procent.

\textbf{7.} Fe/Mg $= 6.26/22.8 = 0.2746$; Ni/Mg $= 0.196/22.8 = 0.00860$.

\textbf{8.} IJzer: $1 - 0.2746/1.907 = 0.856$, 85.6 procent ontbreekt; nikkel:
$1 - 0.00860/0.1110 = 0.923$, 92.3 procent.

\textbf{9.} Iridium is sterker siderofiel dan nikkel, en nikkel sterker dan ijzer:
hoe meer een element metaal verkiest, hoe meer ervan het metaal naar de kern meenam.

\textbf{10.} In de silicaten, als \ce{Fe^2+} in olivijn en pyroxeen (8 procent FeO):
ijzer dat tot FeO geoxideerd is, bindt zich aan zuurstof als een \omterm{def:b2:elements-in-the-earth:goldschmidt}{lithofiel element} en bleef in de
mantel; alleen het metallische ijzer zonk.

\textbf{11.} $(1-x)C_{\mathrm{Fe}} + xc_{\mathrm{Fe}} = r_{\mathrm{Fe}}(1-x)C_{\mathrm{Mg}}$;
delen door $(1-x)c_{\mathrm{Fe}}$ geeft de propositie.

\textbf{12.} $r_{\mathrm{Fe}}C_{\mathrm{Mg}} - C_{\mathrm{Fe}} = 1.907\times22.8 - 6.26 =
37.22$; $y = 37.22/85.5 = 0.4353$, $x = y/(1+y) = 0.303$.

\textbf{13.} $0.1110\times22.8 - 0.196 = 2.335$; $y = 2.335/5.2 = 0.4490$, $x = 0.310$.

\textbf{14.} $y = (85.5\times37.22 + 5.2\times2.335)/(85.5^2 + 5.2^2) =
3194.4/7337.3 = 0.4354$, $x = 0.303$.

\textbf{15.} Elke vergelijking weegt als het kwadraat van haar coëfficiënt: $85.5^2/5.2^2 = 270$
keer zwaarder voor ijzer, het veel overvloedigere element.

\textbf{16.} $37.22/2.335 = 15.9$, tegen $85.5/5.2 = 16.4$ in het model: 3 procent
verschil.

\textbf{17.} $0.303\times\num{5.972e24} = \qty{1.81e24}{kg}$.

\textbf{18.} $(32.5-30.3)/32.5 = 7$ procent lager.

\textbf{19.} $y = 0.325/0.675 = 0.4815$, Fe in de kern $= 37.22/0.4815 = 77.3$ procent,
en Ni $= 2.335/0.4815 = 4.85$ procent: 18 procent lichte elementen, ongeveer het dubbele van de 5 tot
10 procent die de dichtheid toelaat. Fe/Mg van de aarde moet iets hoger zijn dan de CI-waarde,
of haar silicaatmagnesium hoger dan dat van het model.

\textbf{20.} Silicium, zwavel, koolstof, zuurstof en waterstof. Silicium is schaars in de
silicaataarde: zijn op Mg genormaliseerde verhouding, 0.83, ligt onder het refractaire niveau.

\textbf{21.} $1.111\times22.8 - 21.0 = 4.331$; $4.331/0.4354 = 9.9$ procent Si in de
kern.

\textbf{22.} 9.9 tegen 6 procent: de kern bevat misschien meer silicium dan het model, of
een deel van het silicium is nooit geaccreteerd. Silicium condenseert bij \qty{1310}{K}, iets
onder magnesium en ijzer (1336 en \qty{1334}{K}): een deel is misschien in het gas gebleven.

\textbf{23.} Kern: $0.325\times0.5 = \qty{0.1625}{ppm}$ van de massa van de aarde; silicaataarde:
$0.675\times0.001 = \qty{0.000675}{ppm}$. De kern bevat
$0.1625/0.1632 = 99.6$ procent van het goud.

\textbf{24.} De chemie alleen plaatst
\textbf{\qty[text-series-to-math=true]{30.3}{\percent}} van de massa van de aarde in haar kern,
binnen 7 procent van de seismologische 32.5 procent.
\end{solution}
